\documentclass[12pt]{article}

\usepackage[a4paper, total={6in, 8in}]{geometry}
\usepackage{textcomp}

\begin{document}
\setlength{\parindent}{0pt}

\def \t {\textrightarrow}
\def \v {\vspace{1em}}
\def \bi {\begin{itemize}}
\def \ei {\end{itemize}}
\def \s[#1] {\section*{#1}}
\def \ss[#1] {\subsection*{#1}}
\def \sss[#1] {\subsubsection*{#1}}

\s[Pasolini]
Youtube: pier paolo pasolini "L'ultima notte" \t ci sono 3 contributi: 
\bi
  \item di Ninetto, primo suo amante per 12 anni
  \item un altro
  \item l'ultimo, a cui si deve l omicidio
\ei
Lui era un poeta, un regista, un giornalista \t la sua produzione varia da "Ragazzi di Vita" che è importante per la spontaneita, per la primitivita (ribaltata rispetto al fanciullino di pascoli) \\
L opera mostra come la visceralita di questi ragazzi è sintomo di grande umanita \\
Contributo del 2025 sul corriere della sera \\
I ragazzi sono personalita complesse per la societa, ma semplici per l uomo \t lui aderisce all idea di verga di farli parlare nella loro visceralita e nella loro umanita poco urbana \\
Lui è stato accusato di molestie a un ragazzo delle medie, perche insegnava alle medie \\
Lui si allontana poi dal Friuli dove insegnava, e va a Roma \t i ragazzi di vita sono quelli che si trovano nelle spiagge romane, ed è un umanita che non riconosce filtri ed è emblema di semplicita \\
Fa anche il "Decamerone" \t dove mostra tutte le avventure che si sottendono nell opera di Boccacio dei personaggi \\
Mondo dell'industrializzazione, del progresso \t della nuova generazione che si interfaccia al boom economico \\
Mario Soldati si è occupato di lui ? \\
Tutte le sue riflessioni da "scritti corsari" e "Lettere mediterranee" \t sono sferzanti nei confronti della societa \t lui analizza la distorsione dei valori \\
Viene cacciato dal PCI ?? \t ma non si lascia abbattere \t in quest'opera ci sono rimandi a "Sentierio dei nidi di ragno" di Calvino \\
L'intellettuale è utile per segnalare alla politica, mentre la politica rifiuta \t questo è quello che accade \\
Come Gramsci, mette a luce la degenerazione dei valori sociali \\
"Le cenere di Gramsci" \t è stato incarcerato durante il fascismo \\
È un intellettuale che parla alle masse \t questo porta a riflessione di chi è l'intellettuale ora ? \t l'intellettuale adesso? \\
Anche il poeta, chi è? \t attraverso queste voci non c'è una carriera \t è il far parlare l animo \\
Pasolini muore nel 75 \\
E ritrae lo scandalo del mondo borghese \t infatti "Ragazzi di vita" quando esce crea scandalo, ma anche con Moravia \t l'intellettuale di quest epoca è propositivo verso la denuncia, è visibile sulle piattaforme mediatiche \t ma dietro non c'è + un partito a cui appellarsi, quindi non puo darsi carattere politico ??? \\
Da quando Pasolini ha parlato di una borghesia falsa, da quando ha demistificato lo status del matrimonio e ha parlato di un umanita che non ha prezzo \t in tutto questo ha salvato pero la figura della madre \\
Scrive un testo riferito alla madre \t la sua ribellione non è verso i veri valori, ma verso la distorsione dei veri valori \t lui ha sempre avuto rispetto per la madre \\
Aspre critiche nei confronti di Pasolini \\
Fu anche giudicato malato mentalmente \\
Questa epoca non comprende + l intellettuale \t mentre Dante e Petraca no

\v

Si sviluppa la letteratura-industria \t c'è una criticita nel 900 dell industrializzazione, che aveva criticato gia verga \t anche manzoni in realta nel ottavo capitolo, e le poesie di primo 900, e lo stesso Montale in "Trionfo della spazzatura"

\s[Mario Volponi]
Si avvicina al marxismo e scrive "Il memoriale" \t con una rievocazione delle memorie di altri intellettuali \\
Fa analisi critica della spersonalizzazione che porta il lavoro, che fatto da Pirandello con "Quaderni di serafino ?" \t anche in relazione all arte con de chirico \\
Nel mondo dell industria, follia e ragione si contendono il primato \t e non trovano un equilibrio \t ma il mondo moderno vede nell industrializzazione un mondo opprimente \\
Le strategie comunicative e il linguaggio sono quotidani, ma si sottende una precisione scientifica e analitica nell'osservare fenomeni industriali \\
C'è un conveniens nel riferisi ai personaggi \t che tornano da lavoro stanchi con salari bassi \\
I sindacali, rappresentano la irrealizzazone utopico nell equilibrio tra uomo e industria \\
"3 operai di Bernari" e Goffredo Parise \t riferimenti che devono essere ricordati

\v

Alda Merini, che arriva alla follia per ?? \t lei è voce del 900, e apre varco sulla considerazoine della persona \\
Qua il rispetto per una sofferenza fatta parola, come in Calvino con giornata di uno scrutatore \\
Lei denuncia ?? \\
Mario Luzzi \t semplicita espressiva \t e la ricerca della donna riporta a Montale, "giuseppina che non vedo da molto tempo" \t e Sandro Penna, con l'omessualita vissuta e percepita in una lirica \\
Di Luzzi "La notte lava la mente", e "Notizie a giuseppina" \t mentre di Penna "???" \t sono tutte dichiarazioni di dolore \\
Bauman ?? parla di una societa liquida
\end{document}
